\documentclass[aspectratio=169,11pt]{beamer}
\usepackage[T1]{fontenc}
\usepackage{lmodern,amsmath,amssymb,xcolor}
\usetheme{Madrid}
\definecolor{calmTeal}{HTML}{0F766E}
\definecolor{softMint}{HTML}{DFF7F4}
\definecolor{softBlue}{HTML}{EDF6FF}
\definecolor{darkText}{HTML}{1F2937}
\definecolor{mutedText}{HTML}{4B5563}
\setbeamercolor{block body}{bg=softBlue,fg=darkText}
\setbeamercolor{example body}{bg=softMint,fg=darkText}
\newcommand{\qbox}[1]{\begin{beamercolorbox}[rounded=true,wd=\linewidth,sep=0.75em]{block body}\color{darkText}\textbf{Question:} #1\end{beamercolorbox}}
\newenvironment{answerbox}[1]{\begin{beamercolorbox}[rounded=true,wd=\linewidth,sep=0.75em]{example body}\color{darkText}\textbf{#1}\par\medskip}{\end{beamercolorbox}}
\newenvironment{cleanblock}[1]{\begin{beamercolorbox}[rounded=true,wd=\linewidth,sep=0.75em]{block body}\color{darkText}\textbf{#1}\par\medskip}{\end{beamercolorbox}}
\newcommand{\nextanswer}{\vfill\centering\textcolor{mutedText}{\Large Answer on the next slide}}
\begin{document}
% ---------- Problem 2: original statement, part 1 ----------
\begin{frame}{Problem 2 - Problem Statement (1/2)}
\small
\begin{cleanblock}{Problem 2 (Problem 2.32 in the book)}
In a digital communication system, a block of $k$ data bits is mapped into an $n$-bit codeword that typically contains the $k$ information bits as well as $n-k$ redundant bits.

\medskip
This is known as an $(n,k)$ block code. The redundant bits are included to provide error correction capability.
\end{cleanblock}
\end{frame}

% ---------- Problem 2: original statement, part 2 ----------
\begin{frame}{Problem 2 - Problem Statement (2/2)}
\small
\begin{cleanblock}{Problem 2 (continued)}
Suppose that each transmitted bit in our digital communication system is received in error with probability $p$.

\medskip
Furthermore, assume that the decoder is capable of correcting any pattern of $t$ or fewer errors in an $n$-bit block.

\medskip
That is, if $t$ or fewer bits in an $n$-bit block are received in error, then the codeword will be decoded correctly, whereas if more than $t$ errors occur, the decoder will decode the received word incorrectly.

\medskip
Assuming each bit is received in error with probability $p=0.03$, find the probability of decoder error for each of the following codes.
\end{cleanblock}
\end{frame}

% ---------- Explain before solving ----------
\begin{frame}{Problem 2 - Interpreting the Symbols}
\footnotesize
\begin{cleanblock}{Reading the question}
\begin{itemize}
\item $n$: total bits in a codeword.
\item $k$: information bits; $n-k$: redundant bits.
\item $t$: maximum number of correctable bit errors.
\item $p=0.03$: probability of error in one bit.
\end{itemize}
\end{cleanblock}
\begin{answerbox}{Define the random variable}
Let $N$ count bit errors in the $n$ transmitted bits.

Under the independent-bit-error model:
\[
N\sim\operatorname{Binomial}(n,0.03).
\]
\end{answerbox}
\end{frame}

% ---------- Derive general formula ----------
\begin{frame}{Problem 2 - Derive the Error Formula}
\small
\begin{cleanblock}{Step 1: When is decoding correct?}
The decoder corrects $N=0,1,\ldots,t$ errors.
\[
P(\text{correct})=P(N\le t).
\]
\end{cleanblock}
\begin{answerbox}{Step 2: Use the complement rule}
Decoder error occurs when $N>t$:
\[
P_e=P(N>t)=1-P(N\le t).
\]
\end{answerbox}
\end{frame}

\begin{frame}{Problem 2 - General Formula}
\footnotesize
\begin{cleanblock}{Step 3: Use the binomial PMF}
\[
P(N=i)={}^{n}C_i\,p^i(1-p)^{n-i}.
\]
Add the correctable cases:
\[
P(N\le t)=\sum_{i=0}^{t}P(N=i).
\]
\end{cleanblock}
\begin{answerbox}{Final formula}
\[
\boxed{P_e=1-\sum_{i=0}^{t}{}^nC_i\,p^i(1-p)^{n-i}}
\]
Here $p=0.03$ and $1-p=0.97$.
\end{answerbox}
\end{frame}

% ---------- Part a ----------
\begin{frame}{Problem 2a - Question}
\qbox{Find the probability of decoder error for
\[
(n,k)=(7,4),\qquad t=1.
\]}
\nextanswer
\end{frame}

\begin{frame}{Problem 2a - Solution}
\small
\begin{answerbox}{Step 1: Identify correctable cases}
$n=7$, $t=1$: the decoder corrects 0 or 1 error.
\[
P_e=1-P(N=0)-P(N=1).
\]
\end{answerbox}
\begin{answerbox}{Step 2: Substitute the binomial PMF}
\[
\begin{aligned}
P_e=1&-{}^7C_0(0.97)^7\\
     &-{}^7C_1(0.03)(0.97)^6.
\end{aligned}
\]
\[
\boxed{P_e\approx0.01709}
\]
\end{answerbox}
\end{frame}

% ---------- Part b ----------
\begin{frame}{Problem 2b - Question}
\qbox{Find the probability of decoder error for
\[
(n,k)=(15,7),\qquad t=2.
\]}
\nextanswer
\end{frame}

\begin{frame}{Problem 2b - Solution}
\small
\begin{answerbox}{Step 1: Identify correctable cases}
$n=15$, $t=2$: the decoder corrects 0, 1 or 2 errors.
\[
P_e=1-P(N=0)-P(N=1)-P(N=2).
\]
\end{answerbox}
\begin{answerbox}{Step 2: Substitute the binomial PMF}
\[
\begin{aligned}
P_e=1&-{}^{15}C_0(0.97)^{15}\\
     &-{}^{15}C_1(0.03)(0.97)^{14}\\
     &-{}^{15}C_2(0.03)^2(0.97)^{13}.
\end{aligned}
\]
\[
\boxed{P_e\approx0.00937}
\]
\end{answerbox}
\end{frame}

% ---------- Part c ----------
\begin{frame}{Problem 2c - Question}
\qbox{Find the probability of decoder error for
\[
(n,k)=(31,16),\qquad t=3.
\]}
\nextanswer
\end{frame}

\begin{frame}{Problem 2c - Solution}
\footnotesize
\begin{answerbox}{Correctable cases and substitution}
$n=31$, $t=3$: correctable errors are 0, 1, 2, 3.
\[
P_e=1-\sum_{i=0}^{3}P(N=i).
\]
\[
\begin{aligned}
P_e=1&-{}^{31}C_0(0.97)^{31}\\
     &-{}^{31}C_1(0.03)(0.97)^{30}\\
     &-{}^{31}C_2(0.03)^2(0.97)^{29}\\
     &-{}^{31}C_3(0.03)^3(0.97)^{28}.
\end{aligned}
\]
\[
\boxed{P_e\approx0.01335}
\]
\end{answerbox}
\end{frame}
\end{document}
